\clearpage
\section*{Exercise 9: Increasing the Green LED Duration}
\addcontentsline{toc}{section}{Exercise 9: Increasing the Green LED Duration}

\begin{questionbox}{1. Problem Statement}
Use Button 2 to modify the green-light duration and Button 3 to confirm the selected value. The amber duration remains unchanged, and the red duration is recalculated to maintain the timing relationship.
\end{questionbox}

\vspace{0.4cm}
\textbf{2. Green-duration Configuration}

Mode 4 loads \texttt{greenTimeSeconds} into the selected timing value. Button 2 increases the selection on a short press, while the supplied code supports a decrement on a long press. Button 3 commits the selected green duration.

The confirmation logic sets the green value and derives the red value from the unchanged amber value:

\begin{lstlisting}[language=C, caption={Green timing confirmation}, label={lst:exercise9_green_confirm}, captionpos=b]
else if (mode == MODE_SET_GREEN) {
    greenTimeSeconds = selectedTimeSeconds;
    redTimeSeconds =
        (uint8_t)(greenTimeSeconds + yellowTimeSeconds);
}
\end{lstlisting}

The valid green range starts at 1 and ends at \(99 - T_{\mathrm{amber}}\), ensuring that the derived red time remains no greater than 99. For example, with amber = 3, increasing green from 7 to 8 and confirming changes red from 10 to 11 while amber stays at 3.

\textbf{3. Display and LED Feedback}

During Mode 4, the vertical display shows the selected green duration and the horizontal display shows mode number 4. The green logical outputs blink; red and amber outputs are off while the system is in configuration mode. The saved timing is used by the normal traffic sequence after the application returns to Mode 1.

\begin{resultbox}
The implementation keeps the green and amber values as the independently editable values and derives red from their sum. This prevents the configuration from violating \(T_{\mathrm{red}} = T_{\mathrm{green}} + T_{\mathrm{amber}}\).
\end{resultbox}

\vspace{0.2cm}
\vspace{0.15cm}
\noindent\textbf{GitHub source repository:} \href{https://github.com/haunguyen-2006/STM32LAB1}{STM32LAB1}.

% IMAGE SLOT: Add the screenshot described below and save it as ex09_green_duration_mode4.png in the images/ folder.
% If the file is absent, a labeled placeholder is displayed; the document remains compilable.
\begin{figure}[H]
    \centering
    \IfFileExists{images/ex09_green_duration_mode4.png}{%
        \includegraphics[width=0.98\textwidth]{images/ex09_green_duration_mode4.png}%
    }{%
        \fbox{\parbox[c][4.4cm][c]{0.92\textwidth}{%
            \centering\textbf{IMAGE SLOT -- Mode 4 green-duration configuration}\\[6pt]
            Insert a Proteus screenshot of Mode 4 showing the green LEDs blinking, the selected green timing value, and mode number 4.\\[6pt]
            \texttt{images/ex09\_green\_duration\_mode4.png}%
        }}%
    }
    \caption{Mode 4: green-duration configuration and LED feedback}
    \label{fig:exercise9_green_mode_simulation}
\end{figure}


